\documentclass[12pt]{article}
\usepackage[margin=1in]{geometry}
\usepackage{enumitem}
\usepackage{titlesec}
\usepackage{parskip}
\usepackage{graphicx}
\usepackage{hyperref}
\usepackage{fontenc}

\title{\textbf{Moore 1966 RG Notes}\\
\large Moore, Barrington, Jr. \textit{Social Origins of Dictatorship and Democracy: Lord and Peasant in the Making of the Modern World}.\\
Boston: Beacon Press, 1966.}
\author{}
\date{}

\begin{document}

\maketitle

\section{Main Idea}



%--------------------------------------------------
\section{General Notes}
%--------------------------------------------------

\begin{itemize}[leftmargin=*, itemsep=4pt]
    \item Central question: ``We seek the understand the role of the landed upper classes and the peasants in the bourgeois revolutions leading to capitalist democracy, the abortive bourgeois revolutions leading to fascism, and the peasant revolutions leading to communism.''
    \item English liberalization largely intertwined with changing (increasing) value of land. In late 1500s-early 1600s, land values rapidly increased, and lords began to stake claim to additional lands (legally or extrajudiciously) that had previously been communally used by the peasantry. Transition from collective/customary ruling to private ownership, largely driven by the \underline{yeomen}: smaller-scale capitalists who began to commercialize, but were not as wealthy as the ``true'' aristocratic class. The yeomen critically were not direct collaborators with the Crown -- and in fact, began to push back as their wealth grew. The origins of modern Parliament subsequently began as a ``committee of landlords'' to satisfy / enfranchise this group post-Civil War.
    \item Development of commercial agriculture (largely in line with this above change in the value of land and a movement towards exploitation rather than sustenance) is a driving factor in the development of a liberal society: it is largely the primary instigator of private markets which drives broad social change and allows for / needs new institutions (Huntington?)
    \item One core factor in development: whether or not the aristocracy was dependent on the crown/state or on the market. If dependent on the market, it was much harder to fully upend, and they would have a much better chance of securing political enfranchisement. If dependent on the state, much more ripe for more violent responses
    \item Difficult to have peaceful transition to enfranchisement when there is a large and effectively powerful peasant class. If the peasant class can be neutralized in the years leading up to a significant change in agricultural output, then the odds of peace are much higher -- but, if the peasantry is a significant force, their drastically increased output can lead to a need for violent revolution. This is largely in line with Huntington: if economic wellbeing expands rapidly but political institutions have not meaningfully developed, there will be unrest and violence until an equilibrium is reached
    \item ``No bourgeois, no democracy.''
    \item Four pre-conditions for democracy:
    \begin{enumerate}
        \item Development of a balance of power between the crown/central authority and the landed aristocracy that prevents excessive royal absolutism;
        \item A turn toward commercial agriculture that weakens rather than reinforces repressive labor controls;
        \item Weakening or destruction of the landed aristocracy as an independent political class (or its fusion with the bourgeoisie);
        \item prevention of an aristocratic-peasant coalition against the bourgeoisie/commercial classes.
    \end{enumerate}
\end{itemize}


\section{Important Specific Factual Claims}

\begin{itemize}[leftmargin=*, itemsep=4pt]
    \item On English political development in light of the expanded wool trade: ``Under the pressure of circumstances, the medieval notion of judging economic actions according to their contribution to the health of the social organism began to collapse. Men ceased to see the agrarian problem as a question of finding the beset method of supporting people on the land and began to perceive it as the best way of investing capital in the land''
    \item British Civil War strengthened the land-owning class while the French Revolution weakened it: in Britain, this transfer of power to the wealthy laid the framework for future (peaceful) transitions but required the violence of the civil war to set it up. Yet, this empowered the large farming tenants to make jumps in agrarian production, and they ultimately drove this economic transition (capital was not provided by the landlords)
\end{itemize}


\section{My Critiques}

\begin{itemize}
    \item
\end{itemize}


\end{document}
